In this section, we finish the proofs of Theorems \ref{thm:main-cusp}
and \ref{thm:main-norm-length} by
%% Aug 6, 2007.  Put in parameterization
finding a parameterization of the cone deformation for which we may
bound the sum $\sum b_j$.

In order to bound $\sum b_j$, we need more explicit formulas for the
$b_j$.  Recall that $b_j$ was defined in Definition \ref{def:bj} as:
$$b_j = b_{V_j}(\eta_0, \eta_0).$$
Here $\eta_0$ is the real part of the 1--form $\omega_0$ in the $j$-th
component $V_j$ of a tubular neighborhood of the singular locus.  Thus
we begin by revisiting our decomposition of $\omega_0$ in $V_j$.

%%%%%%%%%%%%%%%%%%%%%%%%%%%%%%%%%%%%%%%%%%%%%%%%%%%%%%%%%%%%%%%%%
	
\subsection{Convex combinations of deformations}
	\label{subsec:convex}

Recall that corresponding to an infinitesimal deformation is the
1--form $\omega_0$.  We want $\omega_0$ to have certain desired
properties.

By equation (\ref{eqn:omega_0-decomp}), we can decompose $\omega_0 =
(\omega_0)_j$ in $V_j$ as
$$(\omega_0)_j =u_j\, \omega_m + (x_j+iy_j)\,\omega_\ell,$$
with $u_j$, $x_j$, and $y_j$ in $\RR$, and where $\omega_m$ and
$\omega_\ell$ were computed in \cite{hk:cone-rigid}.

In \cite{hk:univ}, it was shown that if the singular locus has just
one component, then the entire cone deformation can be parameterized
by the square of the cone angle $t=\alpha^2$.  In this case, $u=u_1$
was determined explicitly in \cite{hk:univ}: $u = -1/(4\alpha^2),$
where $\alpha$ is the cone angle at the core of the singular solid
torus $V=V_1$.

When we have multiple components of the singular locus, we may not
necessarily be able to parameterize the cone deformation in this
manner.  However, we do have the following information.

First, for each $j$ in $\{1, \dots, n\}$, locally we have a
deformation given by changing the $j$-th cone angle only and leaving
the others fixed.  This is a deformation with only one component of
the singular locus, so the methods of \cite{hk:univ} apply, and this
deformation can be parameterized by $t=\alpha_j^2$.  Then
\begin{equation*}
(\omega_0)_j = -\frac{1}{4\alpha_j^2}\, \omega_m +
(x_{jj}+iy_{jj})\omega_\ell.
%\label{eqn:omegaj-1}
\end{equation*}

Changing the angle in the $j$-th tube $V_j$ and leaving the others
fixed also affects the $k$-th tube $V_k$.  Since there is no change of
cone angle in this tube, $(\omega_0)_k$ can be expressed as
\begin{equation*}
  (\omega_0)_k = (x_{jk}+iy_{jk})\omega_\ell.
\end{equation*}

Any of these $j$ deformations may be scaled by a factor $s_j\geq 0$,
and we may take non-negative linear combinations.  These correspond to
new local deformations, in which the rates of change of cone angles
vary according to the choice of the $s_j$.  In a neighborhood of the
$j$-th component, write:
\begin{equation}
(\omega_0)_j = -\frac{s_j}{4\alpha_j^2}\,\omega_m + \sum_{k=1}^n (s_k
x_{kj} + i\,s_k y_{kj})\,\omega_\ell.
\label{eqn:omega_0-S}
\end{equation}

Given a point $\bar{s} = (s_1, s_2, \dots, s_n)$, we may use these
equations for $\omega_0$ to compute the forms $\eta_0$ explicitly, and
thus compute the boundary terms $b_j = b_j(\bar{s})$ explicitly.
Hodgson and Kerckhoff did the calculations without the extra $\bar{s}$
(page 382 of \cite{hk:univ}).  When the $\bar{s}$ is put in, we
obtain:
\begin{equation}
\frac{b_j(\bar{s})}{\mbox{Area}(\partial V_j)}= s_j^2\,c +
\left( (\sum_k s_k x_{kj})^2 + 
   (\sum_k s_k y_{kj})^2 \right)\,a +
  s_j(\sum_k s_k x_{kj}) \,b.
\label{eqn:b_j}
\end{equation}
Here the terms $a$, $b$, and $c$ are constants in terms of $R$ which
agree with those on page 383 of \cite{hk:univ}.  However, we are not
immediately concerned with their values.  More important, $a$ is
always negative, $b$ and $c$ always positive.

From the formula (\ref{eqn:b_j}) we obtain the following information.
\begin{enumerate}
  \item $b_j(\bar{s})$ is quadratic in $\bar{s}$.
  \item If any $s_j=0$ for some $\bar{s}$, then $b_j(\bar{s})\leq 0$.
\end{enumerate}

We also know that the sum of all the $b_j$'s is strictly positive,
provided $\bar{s} \neq (0, 0, \dots, 0)$, by equation
(\ref{eqn:b_j>0}), as well as the remark right after that equation.
We will use this information in \S\ref{subsec:Vnonempty}.

%%%%%%%%%%%%%%%%%%%%%%%%%%%%%%%%%%%%%%%%%%%%%%%%%%%%%%%%%%%%%%%%%


\subsection{Selecting local deformations}
  \label{subsec:Vnonempty}

Each choice of $\bar{s} = (s_1, \dots, s_n)$ corresponds to a local
deformation with 1--form $(\omega_0)_j$ expressed in $V_j$ as in
equation (\ref{eqn:omega_0-S}).
In terms of the deformation, varying $\bar{s}$ varies the rates at
which the cone angles are changing (instantaneously).  When
$\bar{s}=(1, \dots, 1)$, all cone angles are changing at the same
rate.  It might be convenient to let $\bar{s}=(1, \dots, 1)$, to
simplify calculations.  However, such a choice may not give us the
bounds on $\sum b_j$ that we need.

In particular, we would like to select $\bar{s}$ such that we may use
certain inequalities from \cite{hk:univ}.  (These inequalities appear
in the proof of Lemma \ref{lemma:b_j-bound} in this paper.)  In order
to use these inequalities, we will need to find $\bar{s}$ such that
each $b_j(\bar{s}) \geq 0$.

The existence of such an $\bar{s}$ has been discovered by
Hodgson and Kerckhoff, but their result is unpublished, so we include
it as Lemma \ref{lemma:Vnonempty} here.

\begin{lemma}[Hodgson--Kerckhoff]  There is an $\bar{s}\neq 0$ for which
$b_j(\bar{s})\geq 0$ for all $j=1, 2, \dots, n.$
\label{lemma:Vnonempty}
\end{lemma}

\begin{proof}
It suffices, by rescaling, to restrict $\bar{s}$ to the simplex
$$T=\{\bar{s}=(s_1, \dots,s_n) | \sum_j s_j = n, s_j \geq 0\}.$$

We first set up some notation.

Let the $j$-th vertex of $T$ be denoted $t_j$.  So $t_j = (0, \dots,
0, n, 0, \dots, 0)$ where the single non-zero value is in the $j$-th
location.  Notice that at $t_j$, $b_i(t_j)\leq 0$ for all $i\neq j$.
Because the sum of the $b_i$ is always strictly positive, we must have
$b_j(t_j)>0$.

Let the face of $T$ opposite $t_j$ be denoted $f_j$.  That is, on
$f_j$ the $j$-th coordinate is always $0$.  Notice that on $f_j$, $b_j
\leq 0$.  We will assume that $b_j < 0$ on $f_j$.  If not, consider
instead level sets of $b_j$ for values $\epsilon > 0$ arbitrarily
close to $0$.  Below, we will be finding points where all the $b_j$
are positive.  Taking a limit will at least give points where the
$b_j$ are nonnegative.

We will also assume level sets $b_j=0$ intersect transversely.  If
not, take an arbitrarily small perturbation of the quadratic functions
(keeping them quadratic).

First we note that the set where $b_j>0$ is an open neighborhood of
$t_j$ in $T$.  This is because on any line from $t_j$ to $f_j$, $b_j$
goes from positive to negative.  Because $b_j(s)$ is quadratic in $s$,
$b_j=0$ on this line in exactly one value.  Thus the level set $b_j=0$
separates $t_j$ from $f_j$.  It has a well--defined positive and
negative side.

Now, we prove the following sublemma by induction.

\begin{lemma}
  Consider the $k-1$ simplex $S_k$ of $T$ where $s_i=0$ for all $i>k$.
  There exists a nonempty open set $B_k \subseteq S_k$ where $b_j>0$
  for all $j\leq k$.  Furthermore, the set $E_k\subseteq S_k$ where
  $b_j=0$ for all $j<k$ is an odd number of points, each on the
  boundary of $B_k$.  
\label{lemma:sub-vnonempty}
\end{lemma}

Notice that in the above statement, when $k=n$, the simplex $S_n$ is
all of $T$.  Then if the statement is true, $B_n$ is an open set on
which all the $b_j$ are positive, so this will prove Lemma
\ref{lemma:Vnonempty}. 

An additional fact we will pick up from the proof is that the set
$E_n$ where $b_j=0$ for $j=1, 2, \dots, n-1$ is an odd number of
points.  (We will use this additional fact in the proof of Lemma
\ref{lemma:choose-S} below.)

\begin{proof}
We first prove Lemma \ref{lemma:sub-vnonempty} when $k=2$.  In this
case, $S_2$ is the 1--simplex from $t_1$ to $t_2$.  At $t_1$, $b_1>0$.
At $t_2$, $b_1 <0$.  Since $b_1$ is quadratic, somewhere on $S_2$ is a
single point $e_2$ where $b_1=0$.  Notice that $b_2(e_2)$ must be
positive.  This is because each $b_j(e_2)\leq 0$ for $j>2$ (since the
corresponding coordinates $s_j$ are all zero).  Since $b_1(e_2)$ is
zero, and the sum $\sum b_i >0$, this forces $b_2>0$.  Then in this
case, $E_2$ contains the single point $e_2$.  $B_2$ is nonempty, with
$e_2$ on its boundary.

\begin{figure}
  \input{figures/findv.pstex_t}
  \caption{Proof of Lemma \ref{lemma:Vnonempty} when $k=3$.}
  \label{fig:findv}
\end{figure}

We will now prove the lemma by induction for $k+1$.  For reference, an
example of the case when $k+1=3$ is illustrated in Figure
\ref{fig:findv}.

$S_{k+1}$ is the $k$ subsimplex of $T$ with vertices $t_1, t_2, \dots,
t_{k+1}$, on which $s_{k+2}= \dots = s_n=0$.  It has as one boundary
face the $k-1$ simplex $S_k$, on which $s_{k+1}= s_{k+2}= \dots
=s_n=0$.  On any other face of $S_{k+1}$, $s_{k+2}=\dots=s_n=0$, and
also $s_j=0$ for some $j<k+1$.

By induction, we are assuming that on the face $S_k$ the set $B_k$,
where $b_j>0$ for all $j\leq k$ is a nonempty open set.  By
continuity, the set $\hat{B}_k$ in $S_{k+1}$ where the same $b_j>0$
must also be nonempty.  %% In its boundary is the 1--manifold $\hat{E_k}$
%% where $b_j=0$ for $j<k$.
The set $\hat{E}_k \subset S_{k+1}$ where $b_j=0$ for $j<k$ is a
$1$--manifold with boundary.  A neighborhood of $E_k$ in $\hat{E}_k$
lies in the boundary of $\hat{B}_k$.

Consider $\hat{E}_k$.  As a 1--manifold, it consists of closed
components and arcs.  The part of its boundary contained in $S_k$ is
the set $E_k$, so consists of an odd number of points.  Hence there
must be an odd number of arc components of $\hat{E}_k$ with exactly
one endpoint in $S_k$.

Consider the other boundary components of $\hat{E_k}$.  They must lie
on faces of $S_{k+1}$.  Recall the set $\{b_j=0\}$ does not intersect
any face $f_j$ where $s_j=0$.  Thus the only possible additional face
of intersection of $\hat{E_k}$ and $S_{k+1}$ is the face on which
$s_k=s_{k+2}=\dots=s_n=0$, spanned by vertices $t_1, \dots, t_{k-1}$
and $t_{k+1}$ (note this face is in $f_k$).  By induction (using the
induction hypothesis with vertices renumbered), there are an odd
number of intersection points of $\hat{E_k}$ on this face of
$S_{k+1}$.

Thus there must be an odd number of arcs of $\hat{E_k}$ running from
$S_k$ to the face $f_k$.  Along any of these arcs, $b_k$ goes from
positive (at a point of $E_k$ on $S_k$) to negative (since $b_k<0$ on
$f_k$).  Thus it must be zero at an odd number of interior points.
Each of these points is in $E_{k+1}$.  We may also have points of
$E_{k+1}$ coming from the intersection of $\{b_k=0\}$ with closed
components of $\hat{E}_k$, or with arc components of $\hat{E}_k$ with
both endpoints of the arc on the same face (either $S_k$ or $f_k$).
In the former case, $\{b_k=0\}$ must intersect any closed component an
even number of times, and in the latter case $b_k$ has the same sign
on both endpoints, so again $\{b_k=0\}$ intersects the arc an even
number of times.  Thus the total number of points in $E_{k+1}$ is odd.

At each point in $E_{k+1}$, we have $b_1=b_2=\dots=b_k=0$, and we have
$b_{k+2}$ through $b_n$ all less than or equal to zero.  Since the sum
of all the $b_j$ is positive, $b_{k+1}$ must be positive.  This
implies $B_{k+1}$ is a nonempty open set, and the points of $E_{k+1}$
are on its boundary.
\end{proof}

This completes the proof of Lemma \ref{lemma:Vnonempty}.
\end{proof}

Now, as mentioned above, the choice of $\bar{s}$ is tied to the rates
at which the cone angles are varying.  Lemma \ref{lemma:Vnonempty}
implies that a desired $\bar{s}$ exists, but gives us no further
information about $\bar{s}$.  In particular, some cone angle may reach
$2\pi$ before the others.  When this happens, the hyperbolic structure
in the corresponding tube about the singular locus is nonsingular.  We
may thus view the deformation from this point on as a deformation on a
cone manifold with one fewer component of the singular locus.  We do
this, obtaining deformations with fewer and fewer singular components,
until each cone angle has reached $2\pi$.



%%%%%%%%%%%%%%%%%%%%%%%%%%%%%%%%%%%%%%%%%%%%%%%%%%%%%%%%%%%%%%%%%

\subsection{Boundary bound}
\label{subsec:bd}

In this subsection, we will show that
%% Aug 6, 2007.  Parameterization
there exists a parameterization of the cone deformation such that the
sum of the terms $b_j$ is of order $n$.  We first need two lemmas.
Lemma \ref{lemma:b_j-bound} gives us an initial bound on each $b_j$.
Lemma \ref{lemma:choose-S} will allow us to improve this bound.  It is
in the proof of Lemma \ref{lemma:b_j-bound}, giving the initial bounds
on the $b_j$, that we use the positivity result of Lemma
\ref{lemma:Vnonempty}.

Finally, in Theorem \ref{thm:bound-bdry-time} at the end of this
subsection, we put these lemmas together to record the main result of
the section: the bound on the sum $\sum b_j$.

\begin{lemma}
  Suppose we have a cone deformation given as a linear combination of
  deformations, as above, where cone angles go from $0$ to $2\pi$
  along the singular locus.  Suppose also that $S=(s_1, \dots, s_n)$
  is chosen such that for each time $t$, $b_j(S)\geq 0$ for all $j$.
  Suppose the tube radius $R$ is bounded below by $R_1 \geq 0.56$ for
  each $t$.  Then
  $$b_j(S) \leq s_j^2 \,C(R_1),$$
  where $C(R_1)$ is a function of $R_1$
  alone which approaches $0$ as $R_1$ approaches infinity.
\label{lemma:b_j-bound}
\end{lemma}

\begin{proof}
  First, since we're considering a fixed $V_j$, we will drop the
  subscript $j$ in our notation throughout the proof (i.e. $V_j=V,$
  $(\eta_0)_j = \eta_0,$ $s_j=s,$ etc).
  
  Modify a calculation (p. 382 and 383 of \cite{hk:univ}) to include
  $s$, and equation (17) of that paper can be written:
$$\frac{b_j(S)}{\mbox{Area}(\partial V)}\; \leq \;\frac{s^2}{8m^4},$$
where
$m = \alpha\sinh^2 R$ is the length of a meridian on $\partial U$.

We will bound the term $\mbox{Area}(\partial V)\, s^2/(8m^4)$.

Let $A=\mbox{Area}(\partial V)$.  Let $h$ denote the height of
$\partial V$, i.e. the length of an arc perpendicular to the
meridian.  So $h=\ell\cosh R$, where $\ell$ is the length of the
singular locus of the solid torus $V$, and $A=mh$.

Consider
$$\Phi = \frac{8m^4}{A} = \frac{8m^3}{h} =
\frac{8\alpha^3\sinh^3R}{\ell\cosh R}.$$
We will bound $\Phi$ below, therefore bounding its reciprocal above.

Let
$$\Upsilon(R) = \frac{3.3957 \tanh R }{2\cosh 2R}.$$
It was shown in
\cite{hk:univ} (see Theorem 4.4, Corollary 5.1, and the comments on
the multicusp case on p. 411), that $\alpha\ell \geq \Upsilon(R)$.

Then
$$\Phi = \frac{8\alpha^3\sinh^3R}{\ell\cosh R} =
\frac{4\alpha^3}{\ell}\,\frac{1}{\Upsilon(R)}\,g(R)$$
where
$$g(R) = \frac{2\Upsilon(R)\sinh^3R}{\cosh R} = 2\Upsilon(R)\tanh R\sinh^2
R = \frac{3.3957 \tanh^4 R}{1+\tanh^2 R}.$$

So $$\Phi \geq \frac{4\alpha^3}{\ell}\,\frac{1}{\alpha\ell}\,g(R) =
4\left(\frac{\alpha}{\ell}\right)^2g(R).$$

Note that $g(R)$ is increasing with $R$.  Hence since $R\geq R_1$,
$g(R)$ is bounded below by $g(R_1)$.  Thus we focus on the term
$\alpha/\ell$ in $\Phi$.

Define $u=\alpha/\ell$.  In our hyperbolic cone deformation, the
initial cone angle $\alpha$ is zero.  As $\alpha$ approaches $0$, $u$
approaches $L_0^2$, where $L_0$ is the normalized length of the
curve which becomes the meridian (i.e. bounds a singular solid disc)
under the cone deformation (see equation (37) of \cite{hk:univ}).  We
will estimate $u=(u_0 + \Delta u) = (L_0^2+\Delta u)$.

Reworking the calculations of \cite{hk:univ} to include the speed $s$,
and taking derivatives with respect to cone angle $\alpha$ rather than
time, we may rewrite Proposition 5.6 of that paper as:
$$\frac{1}{\alpha}\,\frac{du}{d\alpha} \geq -\frac{2(1+\tanh^2 R)}
{3.3957\tanh^3R} = -\frac{2\tanh R}{g(R)},$$
provided that $R\geq
R_1\geq 0.531$.

It is at this step, to use Proposition 5.6 of \cite{hk:univ}, that we
have required $S$ to be chosen so that $b_j(S) > 0$.

Thus
$$\Delta u \geq \int_{\alpha=0}^{\alpha=2\pi} -2\alpha\frac{\tanh
  R}{g(R)} d\alpha \geq
-\frac{\tanh R_1}{g(R_1)}\int_{0}^{2\pi} 2\alpha d\alpha =
-(2\pi)^2\frac{\tanh R_1}{g(R_1)}.$$

So
$$u^2 \geq (L_0^2 - (2\pi)^2\tanh(R_1)/g(R_1))^2,$$
provided the term
$L_0^2 - (2\pi)^2\tanh(R_1)/g(R_1)$ is non--negative.  We can ensure
this quantity is non--negative by ensuring $L_0^2$ is large.  This can
be done following \cite{hk:univ}.  The following result is essentially
equation (47) of that paper.

\begin{equation}
L_0^2 \geq \frac{2(2\pi)^2}{3.3957(1-\tanh R)}
  \exp\left( \int_1^{\tanh R} F(w) dw\right) = I(R),
  \label{eq:Lhat_int}
\end{equation}
provided the tube radius is at least $R$ throughout the deformation.
Here $F(w) = -(1+4w+6w^2+w^4)/((w+1)(1+w^2)^2)$ is an integrable
function.

The right hand side $I(R)$ of (\ref{eq:Lhat_int}) is increasing with
$R$, so it can be bounded below by a constant in terms of $R_1$.
$$L_0^2 \geq I(R_1)$$

Now if $R_1\geq 0.56$, then 
$I(R_1) - 2(2\pi)^2\tanh(R_1)/g(R_1) \geq 0$.

Putting this together,
$$\Phi \geq 4\left(I(R_1)-(2\pi)^2\frac{\tanh R_1}{g(R_1)}\right)^2
g(R_1).$$
The right hand side of the above equation is a constant,
$\tild{C}_0(R_1)$.

Thus we have
$$b_j(S) \leq \frac{A}{8m^4}\,s^2 =
  \frac{1}{\Phi}\,s^2 \leq \frac{s^2}{\tild{C}_0(R_1)} = s^2\,C(R_1).$$
\end{proof}

%%%%%%%%%%%%%%%%%%%%%%%%%%%%%%%%%%%%%%%%%%%%%%%%%%%%%%%%%%%%%%%%%
%% INSERT CHOICE OF POINT S HERE

The result of Lemma \ref{lemma:b_j-bound} gives us a bound on the sum
of the $b_j$ in terms of $\bar{s}$.  If we could choose $\bar{s} = (1,
\dots, 1)$ (all cone angles increasing at the same rate), then Lemma
\ref{lemma:b_j-bound} would give a bound on $\sum b_j$ of order $n$.
However, our choice of $\bar{s}$ comes from Lemma
\ref{lemma:Vnonempty}.  By the proof of that lemma, we may assume only
that $\bar{s}$ is in the simplex $T=\{\bar{s}\, |\, \sum s_j=n,
s_j\geq 0\}$.  A priori, $\bar{s}$ could be close to a vertex of $T$,
say $(n, 0, \dots, 0)$, in which case Lemma \ref{lemma:b_j-bound}
gives only a bound of order $n^2$ on $\sum b_j$.

For our applications, we need the bound of order $n$.  We will obtain
this better estimate by again revisiting the choice of the point
$\bar{s}=(s_1, s_2, \dots, s_n)$.  This is done in Lemma
\ref{lemma:choose-S}.  


\begin{lemma}
  There exists a point $S=(s_1, \dots, s_n)$ in $T$ such that each
  $b_j(S)\geq 0$, $j=1, \dots, n$, and for any index $j$ with $s_j
  \neq \min\{s_1, \dots, s_n\}$, we have $b_j(S)=0$.
\label{lemma:choose-S}
\end{lemma}

\begin{proof}

First we set up some notation.

\begin{notation}
  Let $B = \{\bar{s}\, |\, b_j(\bar{s}) \geq 0, j=1, \dots, n\}$ (so
  the interior of $B$ is the set $B_n$ of Lemma
  \ref{lemma:sub-vnonempty}).  Again we may normalize such that the
  points $\bar{s}$ lie in the simplex $T= \{(s_1,\dots,s_n)| \sum s_i
  = n, s_i\geq 0\}$.
  
  We define some sets in $B$.  For each $j=1, 2, \dots, n$, let $e_j$
  be the set of all points in $B$ such that $b_j>0$ and for $i\neq j$,
  $b_i=0$.  Note for $j=n$, $e_n$ corresponds to the set $E_n$ in
  Lemma \ref{lemma:sub-vnonempty}.  (So in particular, $e_n$ consists
  of an odd number of points.)
  
  Let $b_{\{i,j\}}$ be the set $b_{\{i,j\}} = B\cap (\cap_{k\neq i,j}
  \{b_k=0\})$.  Note $b_{\{i,j\}}$ has dimension $1$, with components with
  boundary in $e_i$ and $e_j$.
  
  In general, let $I$ be a $k$-element subset $I= \{i_1, i_2, \dots,
  i_k\} \subseteq \{1, \dots, n\}$.  Let $b_I = B\cap(\cap_{j\notin I}
  \{b_j=0\})$.  Notice $b_I$ is bounded by components $b_{I_j}$ where
  $j\in I$, and $I_j$ is the $(k-1)$-element set obtained by removing
  $j$ from $I$.
  
  Now we define some sets in the simplex $T = \{\bar{s}| \sum
  s_i=n\}$.  Let $T_i$ be the region in which $s_i\leq s_j$ for all
  indices $j$.  $T_i$ has dimension $n-1$.  Let $T_{\{i,j\}} = T_i
  \cap T_j$.  Thus each $T_i$ is bounded by faces $T_{\{i,j\}}$ where
  $j$ ranges over all elements of $\{1, \dots, n\}-\{i\}$, as well as
  faces on $\partial T$.  Generally,
  let $I$ be a $k$-element subset of $\{1, \dots, n\}$.  Let $T_I =
  \cap_{j\in I} T_j$.  Note the boundary faces of $T_I$ are the faces
  given by $T_{I\cup\{k\}}$, where $k$ ranges over all elements of
  $\{1, \dots, n\} - I$, as well as faces in $\partial T$.  See Figure
  \ref{fig:lemma-fix} for an illustration when $n=3$.
\end{notation}

\begin{figure}
	\input{figures/fix-lemma.pstex_t}
	\caption{Faces $T_I$ illustrated when $n=3$.}
	\label{fig:lemma-fix}
\end{figure}

We will make transversality assumptions that allow us to ignore the
boundary faces of $T_I$ on $\partial T$.  In particular, assume that
$B$ lies in the interior of $T$, so the sets $b_I$ do not meet
$\partial T$.  If this is not the case, we may apply the argument
below to level sets $b_I = \epsilon$, then take the limit as $\epsilon
\to 0$.  Since $B$ is closed, we will obtain the desired result.
Similarly, we will assume that the $b_I$'s and the $T_I$'s all
intersect transversely.

We will use the transversality assumption that $B$ does not meet
$\partial T$ in the arguments below.  In particular, we will be
considering the intersection of subsets of $B$ with boundary faces of
$T_I$.  Since $B$ does not meet $\partial T$, this allows us to ignore
the faces of $T_I$ which lie on $\partial T$ in our argument.

%% If not, we may apply the argument to level
%% sets $b_I=\epsilon$ inside $B$, and take the limit as $\epsilon\to 0$.
%% Since $B$ is closed, we obtain the desired result.

Now, note the lemma will be immediately true if we can find some index
$j$ such that $e_j$ is in the region $T_j$: for on $e_j$, $b_i=0$ for
all $i\neq j$, and in $T_j$, $s_j$ is minimum.  So assume this never
happens.

Generally, if for some $k$-element set $I \subseteq \{1, \dots, n\}$,
$b_I \cap T_I$ is nonempty, we obtain a point which satisfies the
lemma.  So we will assume this also never happens.  We will show that
in this case, the point $(1, 1, \dots, 1)$ is in $B$, thus concluding
the proof.  To prove this, we need another sublemma.

\begin{lemma}
Let $\#(\cdot)$ denote the number of elements of a finite set.
Assuming that %% for 
each %% $k$--element set $I \subseteq \{1, \dots, n\}$,
$b_I \cap T_I$ is empty, 
$$\sum_I \#(b_{I\cup\{n\}} \cap \partial T_I) \equiv \sum_{J}
\#(b_{J\cup\{n\}} \cap \partial T_J) \equiv 1 \mod 2,$$
where the first
sum is over all $(k+1)$-element subsets $I$ of $\{1, \dots, n-1\}$,
and the second sum is over all $k$-element subsets $J$ of $\{1, \dots,
n-1\}$.
\label{lemma:inductive-S}
\end{lemma}

Check that the statement of the lemma makes sense:  Note if $J$ is a
$k$-element set, then $b_{J\cup\{n\}}$ is $k$-dimensional.  $T_J$ is
$(n-k)$ dimensional, so its boundary faces are $(n-k-1)$-dimensional.
Thus the intersections in the statement of the lemma are
$0$-dimensional manifolds, consisting of a finite number of
points.

Also notice that when $k=(n-2)$ in the above lemma, there exists only
one $(k+1)$-element subset $I$, namely $I=\{1, \dots, n-1\}$.  So
$b_{I\cup\{n\}}$ is the whole set $B$.  On the other hand, the
boundary $\partial T_I$ in this case must be the point in the
intersection of all the $T_i$, or the point $(1, \dots, 1)$.  Thus the
first sum counts the number of times the point $(1, \dots, 1)$ appears
in $B$.  Hence showing this number is odd will finish the
proof of Lemma \ref{lemma:choose-S}.


\begin{proof}
We prove Lemma \ref{lemma:inductive-S} by induction on $k$.

We start by showing what essentially is the case $k=0$.  That is, we
show
$$\sum_{i=1}^{n-1} \#(b_{\{i,n\}} \cap \partial T_i) \equiv
\sum_{i=1}^{n-1} \#(e_n \cap T_i) \mod 2.$$

Since we are assuming $e_n \cap T_n = \emptyset$, the second sum in
the above equation is just the number of points of $e_n$ (since $e_n$
is disjoint from each $T_i \cap T_j$ by transversality).  That number
of points is odd by Lemma \ref{lemma:sub-vnonempty}, so this will give
us the basis step of our induction.

Consider $b_{\{i, n\}} \cap T_i$.  This is a 1--manifold, consisting
of closed curves and arcs.  Each arc has two endpoints.  So
$$\sum_{i=1}^{n-1} \#\partial(b_{\{i,n\}} \cap T_i) \equiv 0 \mod 2$$

On the other hand,
$$\sum_{i=1}^{n-1} \#\partial(b_{\{i,n\}} \cap T_i) =
\sum_{i=1}^{n-1} \#(\partial b_{\{i,n\}} \cap T_i) +
\sum_{i=1}^{n-1} \#(b_{\{i,n\}} \cap \partial T_i).$$
So $\sum_{i\neq n} \#(b_{\{i,n\}} \cap \partial T_i)$ has the same
parity as $\sum_{i\neq n} \#(\partial b_{\{i,n\}} \cap T_i)$.

Consider $\partial b_{\{i,n\}} \cap T_i$.  Note $\partial b_{\{i,n\}}$
consists of points in $e_i$ and $e_n$.  By assumption, $e_i$ is never
in $T_i$.  Thus
$$\sum_{i=1}^{n-1} \#(\partial b_{\{i,n\}} \cap T_i)
= \sum_{i=1}^{n-1} \#(e_n \cap T_i) = \#(e_n).$$
So $\sum_{i\neq n} \#(b_{\{i,n\}} \cap \partial T_i)$ is odd.

Our proof for general $k$ follows essentially the same lines as the
proof above.

Suppose $1\leq k \leq n-2$.  We need to show
$$\sum_I \#(b_{I\cup\{n\}} \cap \partial T_I) \equiv \sum_{J}
\#(b_{J\cup\{n\}} \cap \partial T_J) \mod 2,$$
where the first sum is over $(k+1)$-element subsets of $\{1, \dots,
n-1\}$ and the second is over $k$-element subsets.

Fix some $(k+1)$-element subset $I$, and again consider
$b_{I\cup\{n\}} \cap T_I$.  This is a 1--manifold.  Its components are
closed curves and arcs.  Again the number of endpoints of
these arcs is even.  So when we sum over all such $I$, 
$$\sum_I \#\partial(b_{I\cup\{n\}} \cap T_I) \equiv 0 \mod 2.$$

On the other hand, 
$$\sum_I \#\partial(b_{I\cup\{n\}} \cap T_I) =
\sum_I \#(\partial b_{I\cup\{n\}} \cap T_I) +
\sum_I \#(b_{I\cup\{n\}} \cap \partial T_I).$$

Thus we consider $\sum_I \#(\partial b_{I\cup\{n\}} \cap T_I).$ We
want it to have the same parity as $\sum_J \#(b_{J\cup\{n\}} \cap
\partial T_J)$, where $J$ ranges over all $k$-element subsets of $\{1,
\dots, n-1\}$.  We show now that the two sums are actually equal,
which concludes the proof.

Consider $\partial b_{I\cup \{n\}}$.  This consists of $b_I$, as well
as things of the form $b_{{I_j}\cup\{n\}}$ where $I_j$ is the
$k$-element set obtained from $I$ by removing the element $j\in I$.
Since by assumption, $b_I \cap T_I$ is empty, $\partial b_{I\cup\{n\}}
\cap T_I$ is the union of sets $b_{{I_j}\cup\{n\}} \cap T_I$, where
$j$ ranges over all elements of $I$.

Then
\begin{equation}
\sum_I \#(\partial b_{I\cup\{n\}} \cap T_I) = \sum_I \sum_{j\in I}
\#(b_{{I_j}\cup\{n\}} \cap T_I).
\label{eqn:count1}
\end{equation}

We will change our method of counting in (\ref{eqn:count1}) above.
Replace the terms $I_j$ with $k$-element subsets $J$ of $\{1, \dots,
n-1\}$.  We will run through all such $J$.  Then $I$ will run through
the $(k+1)$-element subsets given by $J\cup\{i\}$, where $i$ is some
element of $\{1, \dots, n-1\} - J$.  So the right hand side of
equation (\ref{eqn:count1}) is equal to
$$\sum_J \sum_{i\in\{1,\dots,n-1\}-J} \#(b_{J\cup\{n\}} \cap
T_{J\cup\{i\}}).$$

On the other hand, we know the boundary $\partial T_J$ is equal to the
union of all faces $T_{J\cup\{j\}}$, where $j$ ranges over all
elements of $\{1, \dots, n\}-J$.  The only such face we are missing in
our sum above is $T_{J\cup\{n\}}$.  But by assumption, the
intersection of $b_{J\cup\{n\}}$ with $T_{J\cup\{n\}}$ is trivial.
Thus
$$\sum_J \sum_{i\in\{1,\dots,n-1\}-J} \#(b_{J\cup\{n\}} \cap
T_{J\cup\{i\}}) = \sum_J \#(b_{J\cup\{n\}} \cap \partial T_J).$$
\end{proof}

This concludes the proof of Lemma \ref{lemma:choose-S}.
\end{proof}

%%%%%%%%%%%%%%%%%%%%%%%%%%%%%%%%%%%%%%%%%%%%%%%%%%%%%%%%%%%%%%%%%

Lemmas \ref{lemma:b_j-bound} and \ref{lemma:choose-S} together give
the desired bound on $\sum b_j$:

\begin{theorem}
  Let $X$ be a hyperbolic manifold with $n+k$ cusps connected to a
  manifold $M$ with $k$ cusps by a cone deformation.  (So in
  particular, there are $n$ components of the singular locus, and cone
  angles at each component go from $0$ to $2\pi$.)  Suppose during the
  deformation, the tube radius about the singular locus remains larger
  than $R_1 \geq 0.56$.
%% 	Then the sum of the boundary terms is
%%   bounded:
	Then the deformation can be parameterized so that at each time: $\min
	s_j =1$ and the sum of the boundary terms satisfies
  $$\sum_{j=1}^n b_j \leq n\,C(R_1).$$
  Moreover, we can assume the
  deformation has reached the hyperbolic structure on $M$ by time
  $t=(2\pi)^2$.
  \label{thm:bound-bdry-time}
\end{theorem}

\begin{proof}
  
  By Lemma \ref{lemma:choose-S}, we may find a value $S$ in $T$ such
  that each $b_j(S) \geq 0$ and such that if for any index $j$, $s_j
  \neq \min\{s_1, \dots, s_n\}$, then $b_j(S)=0$.  Denote the minimum
  value of the $s_j$ by $s$: $s=\min\{s_1,\dots, s_n\}$.  Then Lemma
  \ref{lemma:b_j-bound} implies that
\begin{equation}
  \sum_{j=1}^n b_j = \sum_{\{j|b_j\neq 0\}} b_j \leq \sum_{\{j|b_j\neq
  0\}} s^2\,C(R_1) \leq n s^2\,C(R_1).
\label{eqn:sum-b_j}
\end{equation}

First, we need to check that $s$ is not equal to zero.  Suppose not.
That is, suppose $s=0$.  Then by (\ref{eqn:sum-b_j}), $\sum_{j=0}^n
b_j = 0$.  Recall that by equation (\ref{eqn:b_j>0}) and the remark
right after it, the sum of all $b_j$'s equal zero only if the entire
deformation is trivial.  But this is possible only if each $s_j=0$.
Since the point $(0, 0, \dots, 0)$ is not in the simplex $T$, this is
impossible.

Now, recall the terms $s_j$ affect the rate of change of the
deformation.  Thus we may rescale them without changing the tube
radius.  For the point $S$ in $T$, divide each $s_j$ by $s$.
Replace $S$ by this rescaled point, and denote the rescaled
point by the same notation: $S=(s_1, \dots, s_n)$.  So each $s_j\geq
1$.  Also, $s=\min\{s_j\}=1$.

Thus equation (\ref{eqn:sum-b_j}) gives:
$$\sum_{j=1}^n b_j \leq n\,C(R_1).$$
  
Finally, to notice that the deformation is complete by time
$t=(2\pi)^2$ as claimed, recall that if $s_j=1$ throughout the
deformation, then the cone angle is changing at a rate at which
$\alpha_j$ will reach $2\pi$ at time $t=(2\pi)^2$.  (See the
discussion at the beginning of \S\ref{subsec:convex}.)  Since $s_j\geq
1$ always, the rate of change may only be faster.  Thus each cone
angle reaches $2\pi$ by the time $t=(2\pi)^2$.
\end{proof}

When we combine Theorem \ref{thm:bound-bdry-time} with the results of
Theorem \ref{thm:cusp-1} and Theorem \ref{thm:norm-length-1}, we have
completed the proofs of Theorems \ref{thm:main-cusp} and
\ref{thm:main-norm-length}.


